\documentclass[aps,preprint]{revtex4}%
\usepackage{amsfonts}
\usepackage{amsmath}
\usepackage{amssymb}
\usepackage{graphicx}%
\setcounter{MaxMatrixCols}{30}
%TCIDATA{OutputFilter=latex2.dll}
%TCIDATA{Version=4.10.0.2347}
%TCIDATA{CSTFile=revtex4.cst}
%TCIDATA{Created=Thursday, March 13, 2003 08:58:42}
%TCIDATA{LastRevised=Thursday, April 03, 2003 22:03:50}
%TCIDATA{<META NAME="GraphicsSave" CONTENT="32">}
%TCIDATA{<META NAME="DocumentShell" CONTENT="Articles\SW\REVTeX 4">}
\newtheorem{theorem}{Theorem}
\newtheorem{acknowledgement}[theorem]{Acknowledgement}
\newtheorem{algorithm}[theorem]{Algorithm}
\newtheorem{axiom}[theorem]{Axiom}
\newtheorem{claim}[theorem]{Claim}
\newtheorem{conclusion}[theorem]{Conclusion}
\newtheorem{condition}[theorem]{Condition}
\newtheorem{conjecture}[theorem]{Conjecture}
\newtheorem{corollary}[theorem]{Corollary}
\newtheorem{criterion}[theorem]{Criterion}
\newtheorem{definition}[theorem]{Definition}
\newtheorem{example}[theorem]{Example}
\newtheorem{exercise}[theorem]{Exercise}
\newtheorem{lemma}[theorem]{Lemma}
\newtheorem{notation}[theorem]{Notation}
\newtheorem{problem}[theorem]{Problem}
\newtheorem{proposition}[theorem]{Proposition}
\newtheorem{remark}[theorem]{Remark}
\newtheorem{solution}[theorem]{Solution}
\newtheorem{summary}[theorem]{Summary}
\newenvironment{proof}[1][Proof]{\noindent\textbf{#1.} }{\ \rule{0.5em}{0.5em}}
\begin{document}
\preprint{HEP/123-qed}
\title[Short title for running header]{Long title that appears on the title page}
\author{First Author}
\affiliation{My Institution}
\author{Second Author}
\affiliation{My Institution}
\author{Third Author}
\affiliation{Other Institution}
\keywords{one two three}
\pacs{PACS number}

\begin{abstract}
Shell document for REV\TeX{} 4.

\end{abstract}
\volumeyear{year}
\volumenumber{number}
\issuenumber{number}
\eid{identifier}
\date[Date text]{date}
\received[Received text]{date}

\revised[Revised text]{date}

\accepted[Accepted text]{date}

\published[Published text]{date}

\startpage{101}
\endpage{102}
\maketitle
\tableofcontents


\section{Generalized Strong Orthogonality}

The FORDO can be expanded in natural form as%
\begin{equation}
D^{1}=\sum_{1\leq j\leq p}n_{j}\left\vert \varphi_{j}\right\rangle
\left\langle \varphi_{j}\right\vert \label{nat}%
\end{equation}
We can define a an orthogonal direct sum decomposition of the one particle
Hilbert space as
\begin{equation}
\mathcal{H}^{1}=\mathcal{H}_{D^{1}}\oplus\mathcal{H}_{D^{1\bot}}%
\end{equation}
where
\begin{align}
\mathcal{H}_{D^{1}}  &  =\operatorname{linspan}\left\{  \left.  \varphi
_{j}\right\vert 1\leq j\leq p\right\} \nonumber\\
\mathcal{H}_{D^{1\bot}}  &  =\operatorname{linspan}\left\{  \left.  \chi
_{j}\right\vert 1\leq j\leq q\right\}
\end{align}
and subspaces of $\mathcal{H}^{N}$ based on these by%
\begin{align}
\mathcal{H}_{N-\kappa}^{N}  &  =\operatorname{linspan}\left\{  \left.
\left\vert \varphi_{j_{1}}\cdots\varphi_{j_{\kappa}}\chi_{j_{\kappa+1}}%
\cdots\chi_{j_{N}}\right\rangle \right\vert 1\leq j_{1}<\cdots j_{\kappa}\leq
p;1\leq j_{\kappa+1}<\cdots j_{N}\leq q\right\}  ;0\leq\kappa\leq N\nonumber\\
\mathcal{C}_{\lambda}^{N}  &  =%
%TCIMACRO{\dbigoplus \limits_{\lambda\leq\kappa\leq N}}%
%BeginExpansion
{\displaystyle\bigoplus\limits_{\lambda\leq\kappa\leq N}}
%EndExpansion
\mathcal{H}_{\kappa}^{N};\quad0\leq\lambda\leq N
\end{align}
First we consider the case where $D^{N}$ is a pure state
\begin{equation}
D^{N}=\left\vert \Psi\right\rangle \left\langle \Psi\right\vert
\end{equation}
then
\begin{equation}
\left\vert \Psi\right\rangle =\sum_{1\leq j_{1}<\cdots<j_{N}\leq p}%
C_{j_{1}\cdots j_{N}}\left\vert \varphi_{j_{1}}\cdots\varphi_{j_{N}%
}\right\rangle
\end{equation}
and
\begin{equation}
\left\vert \Psi\right\rangle \left\langle \Psi\right\vert =\sum
_{\substack{1\leq j_{1}<\cdots<j_{N}\leq p\\1\leq k_{1}<\cdots<k_{N}\leq
p}}C_{j_{1}\cdots j_{N}}C_{k_{1}\cdots k_{N}}^{\ast}\left\vert \varphi_{j_{1}%
}\cdots\varphi_{j_{N}}\right\rangle \left\langle \varphi_{k_{1}}\cdots
\varphi_{k_{N}}\right\vert
\end{equation}
Let
\begin{equation}
\breve{T}=\breve{I}+\breve{\Lambda}%
\end{equation}
and consider
\begin{equation}
\breve{T}\left\vert \Psi\right\rangle \left\langle \Psi\right\vert \breve
{T}^{\dagger}=\left\vert \Psi\right\rangle \left\langle \Psi\right\vert
+\breve{\Lambda}\left\vert \Psi\right\rangle \left\langle \Psi\right\vert
+\left\vert \Psi\right\rangle \left\langle \Psi\right\vert \breve{\Lambda
}^{\dagger}+\breve{\Lambda}\left\vert \Psi\right\rangle \left\langle
\Psi\right\vert \breve{\Lambda}^{\dagger}%
\end{equation}
where%
\begin{equation}
\breve{\Lambda}=\sum_{\substack{1\leq j_{1}\cdots<j_{N-2}\leq r\\1\leq
k_{1}\cdots<k_{N}\leq p\\1\leq j_{N-1}<j_{N}\leq q}}\lambda_{j_{1}\cdots
j_{N}k_{1}\cdots k_{N}}\left\vert \psi_{j_{1}}\cdots\psi_{j_{N-2}}%
\chi_{j_{N-1}}\chi_{j_{N}}\right\rangle \left\langle \varphi_{k_{1}}%
\cdots\varphi_{k_{N}}\right\vert
\end{equation}
and%
\begin{equation}
\left\{  \left.  \psi_{j}\right\vert 1\leq j\leq r\right\}  =\left\{  \left.
\varphi_{j}\right\vert 1\leq j\leq p\right\}  \cup\left\{  \left.  \chi
_{j}\right\vert 1\leq j\leq q\right\}
\end{equation}
then
\begin{align}
&  \breve{\Lambda}\left\vert \Psi\right\rangle \left\langle \Psi\right\vert
=\\
&  \sum_{\substack{1\leq j_{1}<\cdots<j_{N}\leq p\\1\leq k_{1}<\cdots
<k_{N}\leq p\\1\leq l_{1}<\cdots<l_{N-2}\leq r;1\leq l_{N-1}<l_{N}\leq
q}}\lambda_{l_{1}\cdots l_{N}j_{1}\cdots j_{N}}C_{j_{1}\cdots j_{N}}%
C_{k_{1}\cdots k_{N}}^{\ast}\left\vert \psi_{l_{1}}\cdots\psi_{l_{N-2}}%
\chi_{l_{N-1}}\chi_{l_{N}}\right\rangle \left\langle \varphi_{k_{1}}%
\cdots\varphi_{k_{N}}\right\vert
\end{align}
Let us first consider $\Psi$ to be an IPS\ then
\begin{equation}
\left\vert \Psi\right\rangle =\left\vert \varphi_{1}\cdots\varphi
_{N}\right\rangle
\end{equation}
and
\begin{equation}
\breve{\Lambda}\left\vert \Psi\right\rangle \left\langle \Psi\right\vert
=\sum\lambda_{j_{1}\cdots j_{N}1\cdots N}\left\vert \psi_{j_{1}}\cdots
\psi_{j_{N-2}}\chi_{j_{N-1}}\chi_{j_{N}}\right\rangle \left\langle \varphi
_{1}\cdots\varphi_{N}\right\vert
\end{equation}
unless the coefficients are all imaginary a sufficient condition for this to
contract to zero under $L_{N}^{1}$ is that there are at least two differences
between the bra and the ket i.e. $\breve{\Lambda}$ is of the form
\begin{equation}
\breve{\Lambda}=\sum_{\substack{1\leq j_{1}\cdots<j_{N-2}\leq r\\1\leq
j_{N-1}<j_{N}\leq q}}\lambda_{j_{1}\cdots j_{N}k_{1}\cdots k_{N}}\left\vert
\psi_{j_{1}}\cdots\psi_{j_{N-2}}\chi_{j_{N-1}}\chi_{j_{N}}\right\rangle
\left\langle \varphi_{1}\cdots\varphi_{N}\right\vert .
\end{equation}
However even if the coefficients are imaginary this is not a sufficient
condition for
\begin{equation}
\breve{\Lambda}\left\vert \Psi\right\rangle \left\langle \Psi\right\vert
\breve{\Lambda}^{\dagger}=\sum\lambda_{j_{1}\cdots j_{N}1\cdots N}%
\lambda_{k_{1}\cdots k_{N}1\cdots N}^{\ast}\left\vert \psi_{j_{1}}\cdots
\psi_{j_{N-2}}\chi_{j_{N-1}}\chi_{j_{N}}\right\rangle \left\langle \psi
_{k_{1}}\cdots\psi_{k_{N-2}}\chi_{k_{N-1}}\chi_{k_{N}}\right\vert
\end{equation}
to contract to zero. which can never be the case, thus one is forced to the
conclusion that
\begin{equation}
\breve{\Lambda}\left\vert \Psi\right\rangle \left\langle \Psi\right\vert =0
\end{equation}


$\rfloor$ Then
\begin{align}
Tr_{N-1}\left\{  \breve{\Lambda}\left\vert \Psi\right\rangle \left\langle
\Psi\right\vert \right\}   &  =0\nonumber\\
Tr_{N-1}\left\{  \left\vert \Psi\right\rangle \left\langle \Psi\right\vert
\breve{\Lambda}^{\dagger}\right\}   &  =0
\end{align}
and
\begin{equation}
Tr_{N-1}\left\{  \breve{\Lambda}^{\dagger}\left\vert \Psi\right\rangle
\left\langle \Psi\right\vert \breve{\Lambda}\right\}  =0
\end{equation}
as in particular
\begin{align}
&  \breve{\Lambda}^{\dagger}\left\vert \Psi\right\rangle \left\langle
\Psi\right\vert \breve{\Lambda}=\\
&  \sum C_{j_{1}\cdots j_{N}}C_{k_{1}\cdots k_{N}}^{\ast}\lambda_{l_{1}\cdots
l_{N}m_{1}\cdots m_{N}}\lambda_{t_{1}\cdots t_{N}v_{1}\cdots v_{N}}^{\ast
}\left\vert \psi_{j_{1}}\cdots\psi_{j_{N-2}}\chi_{j_{N-1}}\chi_{j_{N}%
}\right\rangle \left\langle \psi_{k_{1}}\cdots\psi_{k_{N-2}}\chi_{k_{N-1}}%
\chi_{k_{N}}\right\vert \\
&  \qquad\qquad\qquad\qquad\qquad\times\left\vert \varphi_{l_{1}}\cdots
\varphi_{l_{N}}\right\rangle \left\langle \varphi_{m_{1}}\cdots\varphi_{m_{N}%
}\right\vert \left\vert \psi_{t_{1}}\cdots\psi_{t_{N-2}}\chi_{t_{N-1}}%
\chi_{t_{N}}\right\rangle \left\langle \psi_{v_{1}}\cdots\psi_{v_{N-2}}%
\chi_{v_{N-1}}\chi_{v_{N}}\right\vert
\end{align}
Hence the positive transformation
\begin{equation}
\hat{T}\left(  D^{N}\right)  =\breve{T}D^{N}\breve{T}^{\dagger}=\left(
\breve{I}+\breve{\Lambda}\right)  D^{N}\left(  \breve{I}+\breve{\Lambda
}^{\dagger}\right)  =D^{N}+\breve{\Lambda}D^{N}+D^{N}\breve{\Lambda}^{\dagger
}+\breve{\Lambda}D^{N}\breve{\Lambda}^{\dagger}%
\end{equation}
leading to
\begin{equation}
\hat{\Omega}\left(  D^{N}\right)  =\breve{\Lambda}D^{N}+D^{N}\breve{\Lambda
}^{\dagger}+\breve{\Lambda}D^{N}\breve{\Lambda}^{\dagger}%
\end{equation}
where
\begin{equation}
\hat{T}=\hat{I}+\hat{\Omega}%
\end{equation}
gives
\begin{equation}
Tr_{N-1}\left\{  \hat{\Omega}\left(  D^{N}\right)  \right\}  =0
\end{equation}


If $D^{N}$ is an ensemble state
\begin{equation}
D^{N}=\sum_{1\leq j\leq m}\alpha_{j}\left\vert \Psi_{j}\right\rangle
\left\langle \Psi_{j}\right\vert
\end{equation}
where
\begin{equation}
\sum_{1\leq j\leq m}\alpha_{j}=1;\quad0\leq\alpha_{j}\leq1,\quad1\leq j\leq m
\end{equation}
with the same natural general spin orbitals (ngso) as eq. $\left(
\ref{nat}\right)  $ then
\begin{equation}
\left\vert \Psi_{\kappa}\right\rangle \left\langle \Psi_{\kappa}\right\vert
=\sum_{\substack{1\leq j_{1}<\cdots<j_{N}\leq p\\1\leq k_{1}<\cdots<k_{N}\leq
p}}C_{\kappa j_{1}\cdots j_{N}}C_{\kappa k_{1}\cdots k_{N}}^{\ast}\left\vert
\varphi_{j_{1}}\cdots\varphi_{j_{N}}\right\rangle \left\langle \varphi_{k_{1}%
}\cdots\varphi_{k_{N}}\right\vert ;\quad1\leq\kappa\leq m
\end{equation}
and we once again obtain that $Tr_{N-1}\left\{  \hat{\Omega}\left(
D^{N}\right)  \right\}  =0$. Hence we have the following theorem

\begin{theorem}
If $\left\{  \left.  \varphi_{j}\right\vert 1\leq j\leq p\right\}  $ are the
natural general spin orbitals of the $N$-particle state $D^{N}$ associated
with nonzero occupation numbers then $\hat{T}\left(  D^{N}\right)  $ is also
an $N$-particle state (possibly unormalized) that has the \textbf{same }FORDO.
\end{theorem}

\section{FORDO Preserving Transformations}

Let us consider the unitary transformation
\begin{equation}
\hat{U}\left(  D^{N}\right)  =\breve{U}D^{N}\breve{U}^{\dagger}=e^{i\Upsilon
}D^{N}e^{-i\Upsilon}=D^{N}+i\left[  \Upsilon,D^{N}\right]  +\cdots+
\end{equation}
in order that $\hat{U}$ be a FORDO preserving map
\begin{equation}
Tr\left\{  \left\{  \left[  \Upsilon,D^{N}\right]  +\cdots+\right\}
a_{k}^{\dagger}a_{l}\right\}  =0
\end{equation}
We can set%
\begin{equation}
e^{i\Upsilon}=I+\Lambda
\end{equation}
and obtain
\begin{equation}
\left(  \breve{I}+\breve{\Lambda}\right)  D^{N}\left(  \breve{I}%
+\breve{\Lambda}\right)  ^{\dagger}=D^{N}%
\end{equation}
giving that
\begin{equation}
Tr\left\{  \breve{\Lambda}D^{N}+D^{N}\breve{\Lambda}+\breve{\Lambda}%
D^{N}\breve{\Lambda}^{\dagger}\right\}  =0
\end{equation}


\subsection{Real and imaginary parts of an operator}

First one must define the real and imaginary parts of a vector, this concept
is basis dependent i.e. if
\begin{equation}
\left\vert \upsilon\right\rangle =\sum_{j}\upsilon_{j}\left\vert \varphi
_{j}\right\rangle
\end{equation}
then
\begin{align}
\operatorname{Re}_{\varphi}\left(  \left\vert \upsilon\right\rangle \right)
&  =\sum_{j}\operatorname{Re}\upsilon_{j}\left\vert \varphi_{j}\right\rangle
\\
\operatorname{Im}_{\varphi}\left(  \left\vert \upsilon\right\rangle \right)
&  =\sum_{j}\operatorname{Im}\upsilon_{j}\left\vert \varphi_{j}\right\rangle
\end{align}
Any basis formed by real transformations $T$ of the $\left\{  \varphi
_{j}\right\}  $ basis, so that
\begin{align}
T\left\vert \mathbf{\varphi}\right\rangle  &  =\left\vert \mathbf{\varphi
}\right\rangle \mathbb{T}\nonumber\\
\operatorname{Re}\mathbb{T}  &  \mathbb{=T}\nonumber\\
\operatorname{Im}\mathbb{T}  &  \mathbb{=O}%
\end{align}
gives the same definition of real and imaginary part of a vector. A convenient
basis to define the real and imaginary parts wrt is $\left\{  \left\vert
\mathsf{z}\right\rangle \right\}  $ i.e.
\begin{align}
\left\vert \upsilon\right\rangle  &  =\int\left\langle \mathsf{z}%
|\upsilon\right\rangle \left\vert \mathsf{z}\right\rangle d\mathsf{z}%
=\int\upsilon\left(  \mathsf{z}\right)  \left\vert \mathsf{z}\right\rangle
d\mathsf{z}=\int\left\{  \operatorname{Re}\upsilon\left(  \mathsf{z}\right)
+i\operatorname{Im}\upsilon\left(  \mathsf{z}\right)  \right\}  \left\vert
\mathsf{z}\right\rangle d\mathsf{z}\nonumber\\
\operatorname{Re}_{\mathsf{z}}\left(  \left\vert \upsilon\right\rangle
\right)   &  =\int\operatorname{Re}\upsilon\left(  \mathsf{z}\right)
\left\vert \mathsf{z}\right\rangle d\mathsf{z}\nonumber\\
\operatorname{Im}_{\mathsf{z}}\left(  \left\vert \upsilon\right\rangle
\right)   &  =\int\operatorname{Im}\upsilon\left(  \mathsf{z}\right)
\left\vert \mathsf{z}\right\rangle d\mathsf{z}%
\end{align}
Any basis $\left\{  \varphi_{j}\right\}  $ producing the same definition has
the property
\begin{equation}
\varphi_{j}\left(  \mathsf{z}\right)  \in\mathbb{R}%
\end{equation}
i.e. is composed of \emph{real valued} functions. Such bases form an
equivalence class and the set of all bases can be classified into equivalence
classes that are invariant under real linear transformations.

One can then define an operator $X$ to be real (wrt $\left\{  \left\vert
\mathsf{z}\right\rangle \right\}  )$ when
\begin{equation}
X\left(  \left\vert \mathbf{\varphi}\right\rangle \right)  =\left(  \left\vert
\mathbf{\varphi}\right\rangle \right)  \mathbb{X}_{\varphi};\quad
\mathbb{X}_{\varphi}\mathbb{\in R} \label{con3}%
\end{equation}
Such operators map $\operatorname{Re}_{\mathsf{z}}\left(  \left\vert
\upsilon\right\rangle \right)  $ to $\operatorname{Re}_{\mathsf{z}}\left(
\left\vert \tilde{\upsilon}\right\rangle \right)  $ and $\operatorname{Im}%
_{\mathsf{z}}\left(  \left\vert \upsilon\right\rangle \right)  $ to
$\operatorname{Im}_{\mathsf{z}}\left(  \left\vert \tilde{\upsilon
}\right\rangle \right)  $ i.e preserves the real/imaginary structure defined
by $\left\{  \left\vert \mathsf{z}\right\rangle \right\}  $ basis. Any
operator $Z$ can be represented as
\begin{equation}
Z\left(  \left\vert \mathbf{\varphi}\right\rangle \right)  =\left(  \left\vert
\mathbf{\varphi}\right\rangle \right)  \mathbb{Z}_{\varphi}%
\end{equation}
and representation $\mathbb{Z}$ can be uniquely split into its real and
imaginary parts as
\begin{equation}
\mathbb{Z}_{\varphi}\mathbb{=}\operatorname{Re}\mathbb{Z}_{\varphi}%
\mathbb{+}i\operatorname{Im}\mathbb{Z}_{\varphi}%
\end{equation}
where
\begin{align}
\operatorname{Re}\mathbb{Z}_{\varphi}  &  \mathbb{=}\frac{1}{2}\left(
\mathbb{Z}_{\varphi}+\mathbb{Z}_{\varphi}^{\ast}\right) \nonumber\\
\operatorname{Im}\mathbb{Z}_{\varphi}  &  \mathbb{=}\frac{1}{2i}\left(
\mathbb{Z}_{\varphi}-\mathbb{Z}_{\varphi}^{\ast}\right)
\end{align}
leading to the operator definitions
\begin{align}
\operatorname{Re}_{\mathsf{z}}\left(  Z\right)   &  =\left(  \left\vert
\mathbf{\varphi}\right\rangle \right)  \operatorname{Re}\mathbb{Z}_{\varphi
}\left(  \left\vert \mathbf{\varphi}\right\rangle \right)  ^{\dagger
}\nonumber\\
\operatorname{Im}_{\mathsf{z}}\left(  Z\right)   &  =\left(  \left\vert
\mathbf{\varphi}\right\rangle \right)  \operatorname{Im}\mathbb{Z}_{\varphi
}\left(  \left\vert \mathbf{\varphi}\right\rangle \right)  ^{\dagger}%
\end{align}
and
\begin{equation}
Z=\operatorname{Re}_{\mathsf{z}}\left(  Z\right)  +i\operatorname{Im}%
_{\mathsf{z}}\left(  Z\right)  =\left(  \left\vert \mathbf{\varphi
}\right\rangle \right)  \operatorname{Re}\mathbb{Z}_{\varphi}\left(
\left\vert \mathbf{\varphi}\right\rangle \right)  ^{\dagger}+i\left(
\left\vert \mathbf{\varphi}\right\rangle \right)  \operatorname{Im}%
\mathbb{Z}_{\varphi}\left(  \left\vert \mathbf{\varphi}\right\rangle \right)
^{\dagger}%
\end{equation}
One can observe that $\left\{  \operatorname{Re}_{\mathsf{z}}\left(  Z\right)
,\operatorname{Im}_{\mathsf{z}}\left(  Z\right)  \right\}  $ preserve the
real/imaginary structure defined by the $\left\{  \left\vert \mathsf{z}%
\right\rangle \right\}  $ basis. It also allows us to define the complex
conjugate of an operator as
\begin{equation}
Z^{\ast_{\mathsf{z}}}=\left(  \left\vert \mathbf{\varphi}\right\rangle
\right)  \mathbb{Z}_{\varphi}^{\ast}\left(  \left\vert \mathbf{\varphi
}\right\rangle \right)  ^{\dagger}%
\end{equation}
obviously one can define a different complex conjugation for each class of bases.

We can compare this decomposition which is bases dependent (at least bases
class dependent) to the basis independent decomposition into two hermitian
operators
\begin{equation}
Z=Z_{H1}+iZ_{H2}%
\end{equation}
where
\begin{align}
Z_{H1} &  =\frac{1}{2}\left(  Z+Z^{\dagger}\right)  \nonumber\\
Z_{H2} &  =\frac{1}{2i}\left(  Z-Z^{\dagger}\right)
\end{align}
giving that
\begin{align}
\left(  \mathbb{Z}_{\varphi H1}\right)  _{kl} &  =\frac{1}{2}\left(
Z_{kl}+Z_{lk}^{\ast}\right)  \nonumber\\
\left(  \mathbb{Z}_{\varphi H2}\right)  _{kl} &  =\frac{1}{2i}\left(
Z_{kl}-Z_{lk}^{\ast}\right)
\end{align}
which should be compared to
\begin{align}
\left(  \operatorname{Re}\mathbb{Z}_{\varphi}\right)  _{kl} &  =\frac{1}%
{2}\left(  Z_{kl}+Z_{kl}^{\ast}\right)  \nonumber\\
\left(  \operatorname{Im}\mathbb{Z}_{\varphi}\right)  _{kl} &  =\frac{1}%
{2i}\left(  Z_{kl}-Z_{kl}^{\ast}\right)
\end{align}
which are in general quite different. If $Z=Z^{\dagger}$ then $Z_{lk}^{\ast
}=Z_{kl}$ and
\begin{align}
\left(  \operatorname{Re}\mathbb{Z}_{\varphi}\right)  _{kl}  & =\frac{1}%
{2}\left(  Z_{kl}+Z_{lk}\right)  \nonumber\\
\left(  \operatorname{Im}\mathbb{Z}_{\varphi}\right)  _{kl}  & =\frac{1}%
{2i}\left(  Z_{kl}-Z_{lk}\right)
\end{align}


\subsection{Complex Structures}

One can associate two operations with a given complex structure, multipication
by $i$ $\operatorname{complex}$ conjugation $K_{\mathsf{z}}$ i.e.
\begin{align}
&  J_{\mathsf{z}}:\operatorname{Re}_{\mathsf{z}}\left(  \left\vert
\upsilon\right\rangle \right)  +i\operatorname{Im}_{\mathsf{z}}\left(
\left\vert \upsilon\right\rangle \right)  =i\operatorname{Re}_{\mathsf{z}%
}\left(  \left\vert \upsilon\right\rangle \right)  -\operatorname{Im}%
_{\mathsf{z}}\left(  \left\vert \upsilon\right\rangle \right)  \nonumber\\
&  J_{\mathsf{z}}\left(  \operatorname{Re}_{\mathsf{z}}\left(  \left\vert
\upsilon\right\rangle \right)  \right)  =\operatorname{Im}_{\mathsf{z}}\left(
J_{\mathsf{z}}\left(  \left\vert \upsilon\right\rangle \right)  \right)
\nonumber\\
&  J_{\mathsf{z}}\left(  \operatorname{Im}_{\mathsf{z}}\left(  \left\vert
\upsilon\right\rangle \right)  \right)  =-\operatorname{Re}_{\mathsf{z}%
}\left(  J_{\mathsf{z}}\left(  \left\vert \upsilon\right\rangle \right)
\right)
\end{align}
producing
\begin{align}
J_{\mathsf{z}}\circ\operatorname{Re}_{\mathsf{z}} &  =\operatorname{Im}%
_{\mathsf{z}}\circ J_{\mathsf{z}}\nonumber\\
J_{\mathsf{z}}\circ\operatorname{Im}_{\mathsf{z}} &  =-\operatorname{Re}%
_{\mathsf{z}}\circ J_{\mathsf{z}}%
\end{align}
and
\begin{align}
&  K_{\mathsf{z}}:\operatorname{Re}_{\mathsf{z}}\left(  \left\vert
\upsilon\right\rangle \right)  +i\operatorname{Im}_{\mathsf{z}}\left(
\left\vert \upsilon\right\rangle \right)  =\operatorname{Re}_{\mathsf{z}%
}\left(  \left\vert \upsilon\right\rangle \right)  -i\operatorname{Im}%
_{\mathsf{z}}\left(  \left\vert \upsilon\right\rangle \right)  \nonumber\\
&  K_{\mathsf{z}}\left(  \operatorname{Re}_{\mathsf{z}}\left(  \left\vert
\upsilon\right\rangle \right)  \right)  =\operatorname{Re}\left(
K_{\mathsf{z}}\left(  \left\vert \upsilon\right\rangle \right)  \right)
\nonumber\\
&  K_{\mathsf{z}}\left(  \operatorname{Im}_{\mathsf{z}}\left(  \left\vert
\upsilon\right\rangle \right)  \right)  =-\operatorname{Im}\left(
K_{\mathsf{z}}\left(  \left\vert \upsilon\right\rangle \right)  \right)
\end{align}
producing
\begin{align}
K_{\mathsf{z}}\circ\operatorname{Re}_{\mathsf{z}} &  =\operatorname{Re}%
_{\mathsf{z}}\circ K_{\mathsf{z}}\nonumber\\
K_{\mathsf{z}}\circ\operatorname{Im}_{\mathsf{z}} &  =-\operatorname{Im}%
_{\mathsf{z}}\circ K_{\mathsf{z}}%
\end{align}
The representation of $J_{\mathsf{z}}$ on any orthonormal basis of a complex
Hilbert space $\mathcal{H}$ is
\begin{equation}
J_{\mathsf{z}}\left\vert \mathbf{\varphi}\right\rangle =\left\vert
\mathbf{\varphi}\right\rangle i\mathbb{I}%
\end{equation}
so is not basis dependent. Any complex Hilbert space $\mathcal{H}$ with conb
$\left\{  \varphi_{j}\right\}  $ can be expanded as
\begin{equation}
\mathcal{H=H}_{\mathbb{R}}\oplus i\mathcal{H}_{\mathbb{R}}%
\end{equation}
where
\begin{align}
\mathcal{H}_{\mathbb{R}} &  =\mathbb{R-}\operatorname{linspan}\left\{
\varphi_{j}\right\}  \nonumber\\
i\mathcal{H}_{\mathbb{R}} &  =\mathbb{R-}\operatorname{linspan}\left\{
i\varphi_{j}\right\}
\end{align}
then
\begin{equation}
J_{\mathsf{z}}\left(  \left\vert \mathbf{\varphi}\right\rangle ,i\left\vert
\mathbf{\varphi}\right\rangle \right)  =\left(  \left\vert \mathbf{\varphi
}\right\rangle ,i\left\vert \mathbf{\varphi}\right\rangle \right)  \left(
\begin{array}
[c]{cc}%
\mathbb{O} & \mathbb{I}\\
-\mathbb{I} & \mathbb{O}%
\end{array}
\right)
\end{equation}
and
\begin{equation}
K_{\mathsf{z}}\left(  \left\vert \mathbf{\varphi}\right\rangle ,i\left\vert
\mathbf{\varphi}\right\rangle \right)  =\left(  \left\vert \mathbf{\varphi
}\right\rangle ,i\left\vert \mathbf{\varphi}\right\rangle \right)  \left(
\begin{array}
[c]{cc}%
\mathbb{I} & \mathbb{O}\\
\mathbb{O} & -\mathbb{I}%
\end{array}
\right)
\end{equation}
showing that while both $J_{\mathsf{z}}$ and $K_{\mathsf{z}}$ can be
considered as $\mathbb{R-}$ linear transformations of $\mathcal{H}%
_{\mathbb{R}}\oplus i\mathcal{H}_{\mathbb{R}}$ only $J_{\mathsf{z}}$ can be
extended to a $\mathbb{C-}$linear transformation of $\mathcal{H}$. However a
real Hilbert space $\mathcal{H}_{\mathbb{R}},$ with conb $\left\{  \psi
_{j}\right\}  $

\subsection{Strong Orthogonality via a Variational Condition}

Letting
\begin{equation}
\breve{\Lambda}=\operatorname{Re}_{\mathsf{z}}\breve{A}+i\operatorname{Im}%
_{\mathsf{z}}\breve{B}%
\end{equation}
where
\begin{align}
\breve{A} &  =\frac{1}{2}\left(  \breve{\Lambda}+\breve{\Lambda}^{\ast
}\right)  =\breve{A}^{\ast}\nonumber\\
\breve{B} &  =\frac{1}{2i}\left(  \breve{\Lambda}-\breve{\Lambda}^{\ast
}\right)  =\breve{B}^{\ast}%
\end{align}
and $^{\ast}$ is the complex conjugation defined by the basis $\left\{
\left\vert \mathsf{z}\right\rangle \right\}  .$ Without loss of generality we
can always start from a density operator that has a real representation in a
bases belonging to $\left\langle \left\vert \mathsf{z}\right\rangle
\right\rangle $ as  Using this we obtain%
\begin{equation}
Tr_{N-1}\left\{  \breve{A}D^{N}\breve{A}+\breve{B}D^{N}\breve{B}+\breve
{A}D^{N}+D^{N}\breve{A}\right\}  +iTr_{N-1}\left\{  \breve{B}D^{N}-D^{N}%
\breve{B}\right\}  =0
\end{equation}
so that the real part
\begin{equation}
Tr_{N-1}\left\{  \breve{A}D^{N}\breve{A}+\breve{B}D^{N}\breve{B}\right\}
=-Tr_{N-1}\left\{  \left[  \breve{A},D^{N}\right]  _{+}\right\}
\end{equation}
and the imaginary part
\begin{equation}
Tr_{N-1}\left\{  \left[  \breve{B},D^{N}\right]  \right\}  =0
\end{equation}
\emph{as }these relationships \emph{must} also be true in all represenstations
produced by bases $\left\{  \varphi_{j}\right\}  $ that produce the same
complex conjugation as $\left\{  \left\vert \mathsf{z}\right\rangle \right\}
.$  Leading to
\begin{equation}
Tr\left\{  \left[  \breve{B},D^{N}\right]  a_{k}^{\dagger}a_{l}\right\}  =0
\end{equation}
which can be reeexpressed as
\begin{equation}
Tr\left\{  \breve{B}D^{N}a_{k}^{\dagger}a_{l}-a_{k}^{\dagger}a_{l}D^{N}%
\breve{B}\right\}  =Tr\left\{  \left[  a_{k}^{\dagger}a_{l},\breve{B}\right]
D^{N}\right\}  =0\label{con1}%
\end{equation}
Which is equivalent to finding
\begin{equation}
\delta Tr\left\{  e^{i%
%TCIMACRO{\tsum }%
%BeginExpansion
{\textstyle\sum}
%EndExpansion
\eta_{kl}a_{k}^{\dagger}a_{l}}\breve{B}e^{-i%
%TCIMACRO{\tsum }%
%BeginExpansion
{\textstyle\sum}
%EndExpansion
\eta_{kl}a_{k}^{\dagger}a_{l}}D^{N}\right\}  =0
\end{equation}
i.e.%
\begin{equation}
Tr\left\{  B\left(  \eta_{\ast}\right)  D^{N}\right\}  =\operatorname*{stat}%
_{\eta}Tr\left\{  B\left(  \eta\right)  D^{N}\right\}
\end{equation}
where
\begin{align}
\breve{B}\left(  \eta\right)   &  =U_{1}\left(  \eta\right)  \breve{B}%
U_{1}\left(  \eta\right)  ^{\dagger}\nonumber\\
U_{1}\left(  \eta\right)   &  =e^{i%
%TCIMACRO{\tsum }%
%BeginExpansion
{\textstyle\sum}
%EndExpansion
\eta_{kl}a_{k}^{\dagger}a_{l}}%
\end{align}
The "real" hermitian part $\breve{A}$ is determined by $\breve{B}\left(
\eta_{\ast}\right)  $
\begin{align}
&  Tr\left\{  \left\{  \breve{A}D^{N}\breve{A}+\breve{B}\left(  \eta_{\ast
}\right)  D^{N}\breve{B}\left(  \eta_{\ast}\right)  \right\}  a_{k}^{\dagger
}a_{l}\right\}  =-Tr_{N-1}\left\{  \left[  \breve{A},D^{N}\right]  _{+}%
a_{k}^{\dagger}a_{l}\right\}  \nonumber\\
&  \Leftrightarrow Tr\left\{  \left\{  \breve{A}D^{N}\breve{A}+\left[
\breve{A},D^{N}\right]  _{+}\right\}  a_{k}^{\dagger}a_{l}\right\}
=-Tr\left\{  \breve{B}\left(  \eta_{\ast}\right)  D^{N}\breve{B}\left(
\eta_{\ast}\right)  a_{k}^{\dagger}a_{l}\right\}  \label{con2}%
\end{align}


If $D^{N}$ is a pure IPS\ then eq $\left(  \ref{con1}\right)  $ gives
\begin{equation}
\left\langle \Phi_{\ast}|\left[  a_{k}^{\dagger}a_{l},\breve{B}\right]
\Phi_{\ast}\right\rangle =\operatorname*{stat}_{\eta}\left\langle \Phi_{\ast
}|U_{1}\left(  \eta\right)  ^{\dagger}\breve{B}U_{1}\left(  \eta\right)
\Phi_{\ast}\right\rangle =\operatorname*{stat}_{\eta}\left\langle \Phi
_{ref}|U_{1}\left(  \eta\right)  ^{\dagger}\breve{B}U_{1}\left(  \eta\right)
\Phi_{ref}\right\rangle
\end{equation}
where
\begin{equation}
\left\vert \Phi_{\ast}\right\rangle =U_{1}\left(  \eta_{\ast}\right)
\left\vert \Phi_{ref}\right\rangle
\end{equation}
and RHS of eq $\left(  \ref{con2}\right)  $ gives
\begin{align}
\left\langle \Phi_{ref}|\breve{B}\left(  \eta_{\ast}\right)  a_{k}^{\dagger
}a\breve{B}\left(  \eta_{\ast}\right)  \Phi_{ref}\right\rangle  &
=\left\langle \Phi_{ref}|U_{1}\left(  \eta_{\ast}\right)  ^{\dagger}\breve
{B}U_{1}\left(  \eta_{\ast}\right)  a_{k}^{\dagger}a_{l}U_{1}\left(
\eta_{\ast}\right)  ^{\dagger}\breve{B}U_{1}\left(  \eta_{\ast}\right)
\Phi_{ref}\right\rangle \nonumber\\
&  =\left\langle \Phi_{\ast}|\breve{B}U_{1}\left(  \eta_{\ast}\right)
a_{k}^{\dagger}a_{l}U_{1}\left(  \eta_{\ast}\right)  ^{\dagger}\breve{B}%
\Phi_{\ast}\right\rangle \nonumber\\
&  =\left\langle \Phi_{\ast}|\breve{B}U_{1}\left(  \eta_{\ast}\right)
a_{k}^{\dagger}U_{1}\left(  \eta_{\ast}\right)  ^{\dagger}U_{1}\left(
\eta_{\ast}\right)  a_{l}U_{1}\left(  \eta_{\ast}\right)  ^{\dagger}\breve
{B}\Phi_{\ast}\right\rangle \nonumber\\
&  =\left\langle \Phi_{\ast}|\breve{B}b_{k}^{\dagger}b_{l}\breve{B}\Phi_{\ast
}\right\rangle
\end{align}
and the LHS gives
\begin{equation}
-\left\{  \left\langle \Phi_{\ast}|\breve{A}_{\eta_{\ast}}b_{k}^{\dagger}%
b_{l}\breve{A}_{\eta_{\ast}}\Phi_{\ast}\right\rangle +\left\langle \Phi_{\ast
}|\left[  b_{k}^{\dagger}b_{l},\breve{A}_{\eta_{\ast}}\right]  _{+}\Phi_{\ast
}\right\rangle \right\}
\end{equation}


\section{Forms of the Contraction Map.}

The contraction map
\begin{equation}
L_{N}^{1}:\mathcal{B}_{1}\left(  \mathcal{H}^{N}\right)  \rightarrow
\mathcal{B}_{1}\left(  \mathcal{H}^{1}\right)
\end{equation}
can be expressed in following equivalent ways:

\begin{enumerate}
\item
\begin{align}
L_{N}^{1}\left(  D^{N}\right)   &  =%
%TCIMACRO{\dint }%
%BeginExpansion
{\displaystyle\int}
%EndExpansion
D^{1}\left(  \mathsf{z,z}^{\prime}\right)  \left\vert \mathsf{z}\right\rangle
\left\langle \mathsf{z}^{\prime}\right\vert d\mathsf{z}d\mathsf{z}^{\prime
}\nonumber\\
D^{1}\left(  \mathsf{z,z}^{\prime}\right)   &  =%
%TCIMACRO{\dint }%
%BeginExpansion
{\displaystyle\int}
%EndExpansion
D^{1}\left(  \mathsf{z,z}^{N-1};\mathsf{z}^{\prime},\mathsf{z}^{\prime
N-1}\right)  d\mathsf{z}^{N-1}d\mathsf{z}^{\prime N-1}%
\end{align}


\item
\begin{align}
L_{N}^{1}\left(  D^{N}\right)   &  =%
%TCIMACRO{\dsum \limits_{1\leq j,k\leq r}}%
%BeginExpansion
{\displaystyle\sum\limits_{1\leq j,k\leq r}}
%EndExpansion
D_{jk}^{1}\left\vert \psi_{j}\right\rangle \left\langle \psi_{j}\right\vert
\nonumber\\
D_{jk}^{1}  &  =%
%TCIMACRO{\dsum \limits_{1\leq j_{2},\ldots,j_{N}\leq r}}%
%BeginExpansion
{\displaystyle\sum\limits_{1\leq j_{2},\ldots,j_{N}\leq r}}
%EndExpansion
D_{jj_{2}\cdots j_{N}kj_{2}\cdots j_{N}}^{N}%
\end{align}


\item
\begin{align}
D^{1}  &  =%
%TCIMACRO{\dsum \limits_{1\leq l_{2},\ldots,l_{N}\leq r}}%
%BeginExpansion
{\displaystyle\sum\limits_{1\leq l_{2},\ldots,l_{N}\leq r}}
%EndExpansion
\left\langle \psi_{l_{2}}\cdots\psi_{l_{N}}\right\vert \left\{
%TCIMACRO{\dsum \limits_{\substack{1\leq j,j_{2},\ldots,j_{N}\leq r\\1\leq
%k,k_{2},\ldots,k_{N}\leq r}}}%
%BeginExpansion
{\displaystyle\sum\limits_{\substack{1\leq j,j_{2},\ldots,j_{N}\leq r\\1\leq
k,k_{2},\ldots,k_{N}\leq r}}}
%EndExpansion
D_{jj_{2}\cdots j_{N}kj_{2}\cdots j_{N}}^{N}\left\vert \psi_{j}\psi_{j_{2}%
}\cdots\psi_{j_{N}}\right\rangle \left\langle \psi_{k}\psi_{k_{2}}\cdots
\psi_{k_{N}}\right\vert \right\} \nonumber\\
&  \qquad\qquad\qquad\qquad\qquad\qquad\qquad\qquad\qquad\qquad\qquad
\qquad\qquad\qquad\qquad\times\left\vert \psi_{l_{2}}\cdots\psi_{l_{N}%
}\right\rangle \nonumber\\
&  =Tr_{N-1}\left\{  D^{N}\right\}
\end{align}


\item
\begin{equation}
D_{jk}^{1}=Tr\left\{  a\left(  \psi_{j}\right)  ^{\dagger}a\left(  \psi
_{k}\right)  D^{N}\right\}
\end{equation}

\end{enumerate}


\end{document}